% 【本文件写什么】子模块 ipc-event（\subsection 级聚合）
%   - 职责摘要：事件 IPC 占位 crate。
%   - 子域聚合 lib.rs：os/components/wateros-ipc/ipc-event/src/lib.rs
%   - 如何重导出 api-v0、feature 下挂载哪些 impl
%   - 被谁依赖（syscall、vfs、main 启动链等）
%   - 短综述 + 下方 \input 展开 api 与 impl 叶子
% 【事实来源】
%   - os/components/wateros-ipc/ipc-event/
%   - docs/exports/ 中与 wateros-ipc 相关的 public-api / features 章节

\subsection{ipc-event}

\code{ipc-event}为独立占位 crate，未列入\code{wateros-ipc} workspace members，父聚合\code{ipc.tex}虽\code{\input}本节，但根\code{wateros-ipc/Cargo.toml}无\code{dep:event}，\code{ipc::event}命名空间尚不存在。设计意图是预留「事件对象 / 同步」类 IPC 的 crate 边界，供后续 epoll 事件队列或内核句柄语义落地。

当前实现与顶层\code{api-v0}占位相同：仅\code{add}函数用于依赖解析与单测编译。

\begin{rustcode}
/// 占位算术：仅用于依赖解析与单测编译，
/// 不是事件 IPC 契约的一部分。
#[inline]
pub fn add(left : u64, right : u64) -> u64 { left + right }
\end{rustcode}

\paragraph{后续演进。}
正式接入时应定义事件句柄类型、等待/通知契约，并在\code{wateros-ipc/Cargo.toml}增加\code{feature event}与\code{pub mod event}重导出；当前无内核事件队列、无与\syscall{epoll\_wait}对接的对象实现。
